\subsection{AI as a Strategic National Asset}
\label{sec:8.7.6}
AI capability is now national infrastructure, and states treat it accordingly: as a source of economic advantage, military reach, and administrative control over their own populations. Section~\ref{sec:5.4.3}'s point about dual use is the one that governs here, raised to the scale of a state: what decides how a capability is used is who controls it, what law constrains that control, and whether anyone outside the controlling government can check the result.

China's 2017 “New Generation Artificial Intelligence Development Plan” set an explicit national target: becoming the world leader in AI theory, technology, and application by 2030 \autocite{chinastatecouncil2017newgeneration} \autocite{chinastatecouncil2017newgenerationb}. Other governments have made comparable, if less centrally coordinated, commitments, treating AI leadership the way they treat energy security or advanced manufacturing — a capability no state wants to depend on a rival for. The competition is not limited to the two states with the deepest research base: the United Arab Emirates created the world's first cabinet-level minister of state for artificial intelligence in 2017, treating a national AI strategy as worth its own ministry well before most governments had drafted one at all \autocite{uae2017alolama}.

That competition has increasingly focused on a single physical bottleneck: the advanced chips AI systems run on. Three companies sit at the center of the supply chain that produces them — Taiwan's TSMC, which fabricates roughly nine in ten of the most advanced chips in the world, and the Netherlands' ASML, the sole source of the extreme-ultraviolet lithography machines that fabrication requires — which is why a dispute over Taiwan's status is also, mechanically, a dispute over who gets to build the next generation of AI systems at all \autocite{miller2022chipwar}. The United States has restricted, and periodically revised, the export of its most advanced AI accelerators and chip-manufacturing equipment to China since 2022, on the stated ground that unrestricted access would accelerate Chinese military and surveillance capability, and has pressed the Netherlands and Japan, the two other governments whose companies sit at chokepoints in the same supply chain, to restrict their own exports in parallel \autocite{bis2022implementation}. Whatever one thinks of that specific policy, it makes a general point concrete: AI capability depends on physical infrastructure that a small number of states and companies control, which means AI strategic competition is also, underneath the software, a competition over hardware supply chains.

That dependence has a second half the chip story hides. A fabricated chip has to be installed somewhere, powered, and cooled, and that requirement lands on grids, water systems, and zoning authorities inside particular jurisdictions, on a timescale measured in years. It hands a state a lever over AI capability that needs no export-control regime at all: an interconnection queue, a water allocation, a permit. Section~\ref{sec:6.1} takes up who bears the local cost of those decisions. What matters here is that a government's leverage over the systems built in its territory is wider than the export-control debate suggests, and that much of it is exercised by permitting authorities that were not designed with AI in mind and do not think of themselves as regulators of it.

That competitive framing carries a cost. A government racing to field a capability first has a structural reason to defer the harder question of how the capability gets used, against whom, and under what oversight. Section~\ref{sec:6.4.1} catalogs what AI infrastructure looks like once turned toward those uses: mass biometric surveillance, automated censorship at a scale no human review team could match, and the patchwork of systems popularly, and inaccurately, described as a single Chinese “social credit score.” None of it requires different technology from the uses that fill the rest of this book; it requires a government willing to point an existing capability at its own population, and no external constraint stopping it.

States are no longer the only actors this competition runs through. The cost of training the most capable AI models has grown several-fold every year for the better part of a decade, a trend that puts frontier development out of reach for all but a handful of well-capitalized companies and concentrates effective control of the most capable models in whichever government those companies happen to be headquartered under, regardless of how that government's own public AI spending compares to a rival's \autocite{epoch2024costtrend}. A small number of American and Chinese companies now do most of the work this section's own vocabulary describes as a contest between states, which means a government's actual leverage over its own national AI capability depends on how much of that capability sits inside companies it can regulate, tax, or requisition, and how much sits inside companies it can only ask.

The same infrastructure argument runs the other way. A state with functioning rule-of-law institutions, a free press, and courts a citizen can actually use gets a democratic dividend from identical technology: fraud detection that does not double as a loyalty test, disaster response that does not double as population tracking. Stated that way the dividend is a property a state either has or lacks, and I no longer think that is right. Section~\ref{sec:8.7.10} works out the three ways in which having the institutions fails to be sufficient, and the third of them is the one this book had been missing.

Strategic competition over AI is not, by itself, the problem. It is the setting the next three sections' arguments have to work inside, whether the subject is a weapons program, an international standard, or a domestic statute.
